% Proof fragment for integration by the manuscript editor.
% All notation should be aligned with the surrounding appendix at insertion.
\paragraph{The dominant half.}
For $q\ge1/2$, the right half of $P_{n,q}$ has no strict valley. Here the
right half includes the central difference when $n$ is odd. To prove this,
write
\[
 A_n(k)=P_{n,q}(k)-P_{n,1-q}(k)
       =P_{n,q}(k)-P_{n,q}(n-k).
\]
The forward equation gives
\[
 A_{n+1}(k)=(1-r_n(k))A_n(k)+r_n(k-1)A_n(k-1).
\]
Induction shows $A_n(k)\ge0$ for $2k\ge n$: both parents are in the upper
half unless $k$ is the new centre, where $A_{n+1}(k)=0$. Terms outside the
support vanish. Thus $P_{n,q}(n-k)\le P_{n,q}(k)$ on the upper half, and
the central difference of an odd row is nonnegative.

Retain the augmented first-difference row from the preceding proof, but only
at indices $2k+1\ge n$. Its initial row has pattern contained in $(+,-)$.
Indeed its central value is $2q-1\ge0$, and its outer part consists of the
future right boundary values. Those values have pattern $(+,-)$ because
they are the negatives of the left boundary values for bias $1-q\le1/2$.
At each step the positive two-parent recurrence preserves absence of a
negative value followed by a positive one. The only possible parent outside
the previous upper half occurs at the new odd central difference, which is
nonnegative and therefore cannot begin such a pair. Consequently there are
no $i<j$ with $2i+1\ge n$, $\Delta P_{n,q}(i)<0$ and
$\Delta P_{n,q}(j)>0$. A strict valley would contain such two differences.
The span-two local limit transfers this property to $f_{q,a}$ on
$[0,\infty)$. The same local limit transfers reflection dominance:
$f_{q,a}(-x)\le f_{q,a}(x)$ for $q\ge1/2$ and $x\ge0$.

\paragraph{The interval of biases.}
Put $h=h_a$ and $f_q(x)=qh(x)+(1-q)h(-x)$. Conditional unimodality and
$h(-x)\le h(x)$ for $x\ge0$ imply that $h$ is nondecreasing on
$(-\infty,0]$: a strict descent $h(x)>h(y)$ at $x<y\le0$ would give the
strict valley $(x,y,-x)$.

For $q\ge1/2$, $f_q$ has a strict valley if and only if it has a strict
descent $f_q(x)>f_q(y)$ with $x<y\le0$. One direction follows by using
$-x$ as the right witness and reflection dominance. Conversely, a valley
$(x,y,z)$ must have $x<0$, since the right half has no valley. If $y>0$,
the same half-line property applied to $(0,y,z)$ gives
$f_q(0)\le f_q(y)<f_q(x)$, so $(x,0)$ is the required descent.

For fixed $x<y\le0$, the difference $f_q(x)-f_q(y)$ is affine in $q$
and is nonpositive at $q=1$. If it is positive at $q$, it remains positive
at every $r\in[1/2,q]$. Therefore the valley biases in $[1/2,1]$ are
downward closed. They are relatively open, because two strict inequalities
at three fixed points persist under a small change of $q$, and $q=1$ is
excluded by conditional unimodality. Reflection gives a single interval
$(b_a,1-b_a)$, where $0\le b_a\le1/2$ and $b_a=1/2$ denotes the empty
case. Continuity, integrability and positivity turn a strict valley into
two local maxima, by maximizing on compact intervals on its two sides.
Conversely, two local maxima give a strict valley because the density has
no flat interval. The previously proved global bound excludes a third
maximum. Outside the interval the density is unimodal with exactly one
local maximum, including the two bias endpoints.
